\section{Wyniki pomiarów}
\label{sec:wyniki}
\Cref{tab:wyniki,fig:wyniki} przedstawiają tę samą serię danych
z pliku \texttt{data/measurements.csv}. Jest to przykład powtórzeń
w ustalonych warunkach, a nie poziomów przeznaczonych do sumowania.
\begin{table}[htbp]
  \centering
  \caption{Demonstracyjna seria pomiarów.}
  \label{tab:wyniki}
  \pgfplotstabletypeset[
    col sep=comma,
    columns={measurement,level},
    columns/measurement/.style={column name={Pomiar},column type=c},
    columns/level/.style={column name={\multicolumn{1}{c}{$L_p$ (\unit{\decibel})}},
      column type={S[table-format=2.1]},string type},
    every head row/.style={before row=\toprule,after row=\midrule},
    every last row/.style={after row=\bottomrule}
  ]{data/measurements.csv}
\end{table}
\begin{figure}[htbp]
  \centering
  \begin{tikzpicture}
    \begin{axis}[
      width=0.85\linewidth,height=45mm,
      xlabel={Numer pomiaru},ylabel={$L_p$ (\unit{\decibel})},
      xmin=0.5,xmax=5.5,ymin=69.5,ymax=70.5,xtick={1,2,3,4,5},
      grid=major,grid style={hairline},
      tick label style={font=\small},
      /pgf/number format/use comma,
    ]
      \addplot[accent,mark=*,thick] table[
        x=measurement,y=level,col sep=comma]{data/measurements.csv};
    \end{axis}
  \end{tikzpicture}
  \caption{Rozrzut demonstracyjnych wyników kolejnych powtórzeń.}
  \label{fig:wyniki}
\end{figure}
